
\documentclass[aspectratio=169]{beamer}
\usetheme{Madrid}
\usecolortheme{default}
\usepackage[utf8]{inputenc}
\usepackage[T1]{fontenc}
\usepackage[albanian,english]{babel}
\usepackage{amsmath,amssymb,bm}
\usepackage{graphicx}
\usepackage{booktabs}
\usepackage{siunitx}
\usepackage{ifthen}

\title{Zgjidhja analitike Buckley--Leverett për \textit{polymer flooding}\\dhe lidhja me modelimin e ujërave spitalore}
\author{Klaudio Peqini}
\institute{Material për diskutim me profesorin}
\date{\today}

\newcommand{\safeinclude}[2][0.92\textwidth]{%
  \IfFileExists{#2}{\includegraphics[width=#1,height=0.72\textheight,keepaspectratio]{#2}}%
  {\fbox{\parbox[c][0.45\textheight][c]{0.75\textwidth}{\centering Figura \texttt{#2} nuk u gjet.\\Vendoseni në të njëjtin direktor me këtë skedar ose korrigjoni rrugën.}}}%
}

\begin{document}

\begin{frame}
  \titlepage
\end{frame}

\begin{frame}{Motivimi shkencor}
\begin{itemize}
  \item Profesori kërkoi implementimin e zgjidhjes analitike Buckley--Leverett për \textit{polymer flooding} në Python.
  \item Synimi kryesor është të përshkruhet zhvendosja 1D e naftës nga uji/polimieri në medium poroz, duke përdorur:
  \begin{itemize}
    \item metodën e karakteristikave;
    \item kushtin e tangentës së Welge për formimin e goditjes (\textit{shock});
    \item efektin e rritjes së viskozitetit të fazës ujore për shkak të polimerit.
  \end{itemize}
  \item Paralelisht, u kërkua që ky model të lidhet konceptualisht me paketën ekzistuese të ujërave spitalore, ku janë studiuar \textit{frontet e breakthrough}-ut dhe transporti i ndotësve.
\end{itemize}
\end{frame}

\begin{frame}{Hipotezat e modelit Buckley--Leverett}
\begin{block}{Supozimet bazë}
\begin{itemize}
  \item rrjedhje 1D, e pakompresueshme;
  \item shpejtësi totale Darcy konstante;
  \item pa shpërhapje, pa adsorbim, pa degradim;
  \item viskozitet konstant efektiv i ujit me polimer, $\mu_w^p > \mu_w$;
  \item permeabilitete relative të tipit Corey.
\end{itemize}
\end{block}

\begin{equation*}
\phi \frac{\partial S_w}{\partial t} + \frac{\partial}{\partial x}\left(u f_w(S_w)\right)=0
\end{equation*}

ku:
\[
f_w(S_w)=\frac{\lambda_w}{\lambda_w+\lambda_o}, \qquad
\lambda_w=\frac{k_{rw}}{\mu_w^{(*)}}, \qquad
\lambda_o=\frac{k_{ro}}{\mu_o}.
\]
\end{frame}

\begin{frame}{Permeabilitetet relative dhe rrjedha fraksionare}
Saturimi efektiv:
\[
S_e=\frac{S_w-S_{wi}}{1-S_{wi}-S_{or}}.
\]

Ligjet Corey:
\[
k_{rw}=k_{rw0} S_e^{n_w}, \qquad
k_{ro}=k_{ro0}(1-S_e)^{n_o}.
\]

\begin{itemize}
  \item Rritja e viskozitetit të ujit nga polimeri ul mobilitetin e fazës ujore.
  \item Kjo ndryshon formën e kurbës $f_w(S_w)$.
  \item Si pasojë vonohet \textit{breakthrough}-u i ujit dhe përmirësohet rikuperimi.
\end{itemize}
\end{frame}

\begin{frame}{Metoda e karakteristikave dhe formimi i \textit{shock}-ut}
Nga PDE-ja rezulton shpejtësia karakteristike:
\[
\frac{x}{t}=\frac{u}{\phi}\frac{df_w}{dS_w}.
\]

\begin{itemize}
  \item Çdo vlerë e saturimit lëviz me shpejtësi të ndryshme.
  \item Për shkak të jolinearitetit, karakteristikat mund të kryqëzohen.
  \item Në këtë pikë zgjidhja klasike prishet dhe formohet një front diskret (\textit{shock}).
\end{itemize}

Kushti i Welge:
\[
\left.\frac{df_w}{dS_w}\right|_{S_{wf}}
=
\frac{f_w(S_{wf})-f_w(S_{wi})}{S_{wf}-S_{wi}}.
\]

Shpejtësia dimensionless e frontit:
\[
v_D^{shock}=\frac{f_w(S_{wf})-f_w(S_{wi})}{S_{wf}-S_{wi}}.
\]
\end{frame}

\begin{frame}{Çfarë u korrigjua në kodin fillestar}
\begin{itemize}
  \item U rikthye forma standarde e permeabiliteteve relative Corey; kodi fillestar i normalizonte gabimisht me shumën e dy fazave.
  \item U zëvendësua identifikimi jo i qëndrueshëm i \textit{shock}-ut me kriterin e vërtetë të tangentës së Welge.
  \item U eliminua gabimi i indeksimit të tipit \verb|fw[S_wf]|, ku \verb|S_wf| ishte vlerë reale dhe jo indeks.
  \item U rindërtua profili $S_w(x,t)$ duke inveruar në mënyrë konsistente lidhjen karakteristike.
  \item U bë i qëndrueshëm llogaritja e \textit{water cut}, rikuperimit kumulativ dhe PVI në \textit{breakthrough}.
\end{itemize}
\end{frame}

\begin{frame}{Krahasimi i kurbave të rrjedhës fraksionare}
\centering
\safeinclude{fractional_flow_comparison.png}
\end{frame}

\begin{frame}{Interpretimi fizik i kurbës $f_w(S_w)$}
\begin{itemize}
  \item Tangjentja e Welge përcakton saturimin kritik $S_{wf}$ ku lind fronti i mprehtë.
  \item Në rastin me polimer, kurba e rrjedhës fraksionare zhvendoset në mënyrë të tillë që:
  \begin{itemize}
    \item \textit{breakthrough}-u vonohet;
    \item faktori i rikuperimit në të njëjtën PVI rritet;
    \item ulja e raportit të mobilitetit përmirëson stabilitetin e zhvendosjes.
  \end{itemize}
  \item Ky është pikërisht roli fizik i polimerit në \textit{enhanced oil recovery}.
\end{itemize}
\end{frame}

\begin{frame}{Profilet e saturimit}
\centering
\safeinclude{saturation_profiles_comparison.png}
\end{frame}

\begin{frame}{Kurba të prodhimit: breakthrough dhe rikuperimi}
\centering
\safeinclude{production_curves_comparison.png}
\end{frame}

\begin{frame}{Rezultatet kryesore numerike}
\begin{itemize}
  \item Paketa e re prodhon automatikisht tabela:
  \begin{itemize}
    \item \verb|polymer_summary.csv|
    \item \verb|fractional_flow_table.csv|
    \item \verb|polymer_production_curves.csv|
    \item \verb|polymer_profiles.csv|
  \end{itemize}
  \item Daljet kryesore janë:
  \begin{itemize}
    \item saturimi i \textit{shock}-ut $S_{wf}$;
    \item PVI në \textit{breakthrough};
    \item \textit{water cut} në prodhues;
    \item faktori i rikuperimit $RF$.
  \end{itemize}
  \item Këto rezultate janë të përshtatshme për interpretim fizik, raporte teknike dhe grafikë për publikim/prezantim.
\end{itemize}
\end{frame}

\begin{frame}{Lidhja direkte me modelin ekzistues të ujërave spitalore}
\begin{block}{Ideja qendrore}
Dy modelet i përkasin së njëjtës familje matematikore: \textbf{ligje ruajtjeje me fronte të mprehta dhe breakthrough}.
\end{block}

\begin{center}
\begin{tabular}{ll}
\toprule
Buckley--Leverett & Modeli i ujërave spitalore \\
\midrule
front saturimi $S_w$ & front përqendrimi / breakthrough \\
\textit{water cut} në prodhues & breakthrough i CBZ/DCF në dalje \\
rikuperim kumulativ i naftës & heqje kumulative e ndotësve \\
karakteristika dhe shock & front jolinear i transportit \\
\bottomrule
\end{tabular}
\end{center}

\begin{itemize}
  \item Në paketën e re u shtua edhe një skript lidhës (\verb|run_transport_bridge_example.py|).
\end{itemize}
\end{frame}

\begin{frame}{Figura lidhëse: transporti në medium poroz kundrejt breakthrough-ut të ndotësve}
\centering
\safeinclude{transport_bridge.png}
\end{frame}

\begin{frame}{Pse kjo lidhje është e rëndësishme}
\begin{itemize}
  \item Tregon se puna e kryer më parë për ujërat spitalore nuk është e shkëputur nga detyra e re.
  \item Përkundrazi, të dyja problemet:
  \begin{itemize}
    \item përshkruajnë propagimin e fronteve;
    \item kërkojnë interpretim të \textit{breakthrough}-ut;
    \item analizohen përmes madhësive pa dimension si PVI ose kohë e reduktuar.
  \end{itemize}
  \item Kjo e vendos detyrën e re në një kuadër më të gjerë të modelimit të transportit jolinear.
\end{itemize}
\end{frame}

\begin{frame}{Përmirësime të mundshme të mëtejshme}
\begin{itemize}
  \item Shtimi i adsorbimit të polimerit dhe vonesës (\textit{retardation}).
  \item Shtimi i shpërhapjes për krahasim analitik--numerik.
  \item Kalimi nga viskozitet konstant në reologji më realiste të polimerit.
  \item Futja e \textit{slug injection} dhe \textit{chase water}.
  \item Krahasimi me zgjidhje plotësisht numerike të tipit finite-volume.
\end{itemize}
\end{frame}

\begin{frame}{Përfundime}
\begin{itemize}
  \item U ndërtua një implementim i pastër dhe fizikisht i qëndrueshëm i modelit Buckley--Leverett për \textit{polymer flooding}.
  \item U prodhuan figura të standardit të lartë për:
  \begin{itemize}
    \item rrjedhën fraksionare,
    \item profilet e saturimit,
    \item \textit{water cut} dhe rikuperimin.
  \end{itemize}
  \item U vendos lidhja e drejtpërdrejtë me paketën ekzistuese të ujërave spitalore përmes logjikës së përbashkët të transportit dhe \textit{breakthrough}-ut.
  \item Rezultati final është një paketë Python më e fortë si nga ana shkencore, ashtu edhe nga ana e prezantimit.
\end{itemize}
\end{frame}

\end{document}
